\DocumentMetadata{
  pdfversion=2.0,
  pdfstandard=ua-2,
  lang=en-US,
  tagging=on,
  tagging-setup={math/setup=mathml-SE,tabsorder=structure}
}
\documentclass[11pt,letterpaper]{article}
\usepackage[margin=0.68in,includefoot]{geometry}
\usepackage{newcomputermodern}
\usepackage{amsmath,amscd,mathtools}
\providecommand{\square}{\mdlgwhtsquare}
\usepackage[hidelinks]{hyperref}
\hypersetup{
  pdftitle={Analysis PhD Exam, May 2015},
  pdfauthor={University of Florida Department of Mathematics},
  pdfsubject={Graduate mathematics examination},
  pdfdisplaydoctitle=true,
  bookmarksopen=true
}
\setcounter{secnumdepth}{0}
\setlength{\parindent}{0pt}
\setlength{\parskip}{0.28em}
\setlength{\emergencystretch}{3em}
\setlength{\leftmargini}{2.15em}
\setlength{\leftmarginii}{2.65em}
\setlength{\leftmarginiii}{3em}
\renewcommand{\labelenumi}{\textbf{\theenumi.}}
\pagestyle{empty}
\raggedbottom
\allowdisplaybreaks
\newenvironment{examproblems}[1]{%
  \begin{enumerate}%
  \setcounter{enumi}{#1}%
  \setlength{\topsep}{0.35em}%
  \setlength{\partopsep}{0pt}%
  \setlength{\itemsep}{0.48em}%
  \setlength{\parsep}{0.18em}%
}{\end{enumerate}}
\newenvironment{examsubparts}{%
  \begin{enumerate}%
  \setlength{\topsep}{0.18em}%
  \setlength{\partopsep}{0pt}%
  \setlength{\itemsep}{0.16em}%
  \setlength{\parsep}{0.1em}%
}{\end{enumerate}}
\newenvironment{examinstructions}{%
  \par\begingroup\itshape
}{%
  \par\endgroup\vspace{0.18em}
}
\makeatletter
\renewcommand\section{\@startsection{section}{1}{0pt}%
  {-1.25ex plus -.25ex minus -.1ex}%
  {0.55ex plus .1ex}%
  {\normalfont\large\bfseries}}
\renewcommand{\@maketitle}{%
  \newpage\null\vskip -1.2em%
  {\footnotesize\bfseries
    \href{https://www.ufl.edu/}{University of Florida}, \href{https://math.ufl.edu/}{Department of Mathematics}\par}%
  \vskip 0.18em%
  \tagstructbegin{tag=Title}%
    {\LARGE\bfseries\href{https://gma.math.ufl.edu/exams/analysis-phd/}{Analysis PhD Exam}, May 2015\par}%
  \tagstructend%
  \vskip 0.35em%
  \tagmcbegin{artifact}%
    \hrule height 1pt%
  \tagmcend%
  \vskip 0.55em%
}
\makeatother
\title{Analysis PhD Exam, May 2015}
\author{}
\date{}
\begin{document}
\maketitle
\thispagestyle{empty}
\begin{examinstructions}
Do six of eight. Write solutions in a neat and logical fashion, giving complete reasons for all steps.
\end{examinstructions}
\begin{examproblems}{0}
\item
Let~\(X\) and~\(Y\) be topological spaces. Prove, if \(f:X\to Y\) is continuous, then~\(f\) is \(\mathcal{B}_X\)-\(\mathcal{B}_Y\) measurable (that is, measurable with respect to the Borel~\(\sigma\)-algebras on~\(X\) and~\(Y\) respectively).
\item
(Short Answer.) Do two of three.
\begin{examsubparts}
\item[\textnormal{(i)}]
Give an example, if possible, of a closed subset~\(F\) of~\(\mathbb{R}\) of positive measure which contains no nontrivial open interval.
\item[\textnormal{(ii)}]
Give an example, if possible, which shows that the~\(\sigma\)-finite hypothesis is needed in Tonelli’s Theorem.
\item[\textnormal{(iii)}]
Determine the limit of the sequence, 
\[
\int_1^\infty \frac{\sin(nx)}{x^n}\,dx.
\]
\end{examsubparts}
\item
Let \((X,\mathcal{M},\mu)\) be a measure space and \((f_n)_{n=1}^\infty\) be a sequence from \(L^1(\mu)\). Prove, if there is a \(g\in L^1(\mu)\) such that \(g(x)\geq |f_n(x)|\) for all~\(n\) and \(x\in X\), then 
\[
\lim_{\delta\to 0}\sup_n\int_{\{|f_n|<\delta\}}|f_n|\,d\mu=0.
\]
\item
Prove, if \(E\subset\mathbb{R}\) is a Lebesgue measurable set of positive measure, then \(E-E=\{x-y:x,y\in E\}\) contains a nontrivial open interval.
\item
Let \((X,\mathcal{M},\mu)\) be a measure space and \(1\leq p<q<\infty\). Prove, \(L^q(\mu)\subset L^p(\mu)\) if and only if there is a (real) constant~\(K\) such that for each \(E\in\mathcal{M}\) either \(\mu(E)\leq K\) or \(\mu(E)=\infty\).
\item
Prove, if~\(X\) is a Banach space and~\(M\) is a finite dimensional subspace of~\(X\), then there is a closed subspace~\(N\) of~\(X\) such that \(M\cap N=(0)\) and \(M+N=X\).
\item
Suppose \(\|\cdot\|\) and \(\|\cdot\|_*\) are both norms on a vector space~\(X\) and \(\|x\|\leq\|x\|_*\) for all \(x\in X\). Show, if both \((X,\|\cdot\|)\) and \((X,\|\cdot\|_*)\) are Banach spaces, then these norms are equivalent.
\item
(Short Answer. Do two of three)
\begin{examsubparts}
\item[\textnormal{(a)}]
Explain why \(c_{00}\) can not be a Banach space.
\item[\textnormal{(b)}]
Explain why there is no bounded one-one onto linear map from \(c_0\) to \(\ell^\infty\).
\item[\textnormal{(c)}]
Explain what is meant by the Fourier transform of a function \(f\in L^2(\mathbb{R})\).
\end{examsubparts}
\end{examproblems}
\end{document}
